\section*{第 7 章}
\addcontentsline{toc}{section}{第 7 章}

\exercise{习题 7.1}
\textbf{题目：}
设一信源由 6 个不同的独立符号组成：
\[
x_1,x_2,x_3,x_4,x_5,x_6
\]
其概率分别为
\[
\frac12,\ \frac14,\ \frac1{32},\ \frac18,\ \frac1{32},\ \frac1{16}.
\]
求：
\begin{enumerate}
\item 信源符号熵 $H(X)$；
\item 若信源每秒发送 $1000$ 个符号，求每秒传送的信息量；
\item 若各符号等概出现，求信源最大熵 $H_{\max}(X)$。
\end{enumerate}

\par\textbf{解：}
\[
H(X)=-\sum_{i=1}^6 p_i\log_2 p_i.
\]
代入各概率，
\[
H(X)=\frac12\cdot 1+\frac14\cdot 2+\frac1{32}\cdot 5+\frac18\cdot 3+\frac1{32}\cdot 5+\frac1{16}\cdot 4
=\frac{31}{16}=1.9375\ \text{bit}.
\]

若每秒发送 $1000$ 个符号，则信息量速率为
\[
R=1000\times H(X)=1937.5\ \text{bit/s}.
\]

若 6 个符号等概出现，则最大熵为
\[
H_{\max}(X)=\log_2 6\approx 2.585\ \text{bit}.
\]

\medskip
\exercise{习题 7.2}
\textbf{题目：}
已知两个二进制随机变量 $X,Y$ 的联合分布满足
\[
P(X=0,Y=0)=P(X=0,Y=1)=P(X=1,Y=1)=\frac13,
\]
其余联合概率为 0。求
\[
H(X),\ H(Y),\ H(X|Y),\ H(Y|X),\ H(X,Y).
\]

\par\textbf{解：}
联合熵为
\[
H(X,Y)=-3\cdot \frac13\log_2\frac13=\log_2 3\approx 1.585\ \text{bit}.
\]

边缘分布为
\[
P(X=0)=\frac23,\qquad P(X=1)=\frac13,
\]
\[
P(Y=0)=\frac13,\qquad P(Y=1)=\frac23.
\]
因此
\[
H(X)=H(Y)=-\frac23\log_2\frac23-\frac13\log_2\frac13\approx 0.918\ \text{bit}.
\]

条件熵满足
\[
H(X|Y)=H(X,Y)-H(Y)=\log_2 3-0.918\approx 0.667\ \text{bit},
\]
\[
H(Y|X)=H(X,Y)-H(X)\approx 0.667\ \text{bit}.
\]

\medskip
\exercise{习题 7.3}
\textbf{题目：}
已知联合事件概率表如下：
\[
\begin{array}{c|ccc}
P(A_i,B_j) & B_1 & B_2 & B_3\\
\hline
A_1 & 0.10 & 0.08 & 0.13\\
A_2 & 0.05 & 0.03 & 0.09\\
A_3 & 0.05 & 0.12 & 0.14\\
A_4 & 0.11 & 0.04 & 0.06
\end{array}
\]
求：
\begin{enumerate}
\item $P(A_i)$ 与 $P(B_j)$；
\item $H(A),H(B),H(A,B)$；
\item 互信息 $I(A;B)$。
\end{enumerate}

\par\textbf{解：}
各行和各列求和得
\[
P(A_1)=0.31,\quad P(A_2)=0.17,\quad P(A_3)=0.31,\quad P(A_4)=0.21,
\]
\[
P(B_1)=0.31,\quad P(B_2)=0.27,\quad P(B_3)=0.42.
\]

联合熵为
\[
H(A,B)=-\sum_{i,j}P(A_i,B_j)\log_2 P(A_i,B_j)\approx 3.447\ \text{bit}.
\]

边缘熵分别为
\[
H(A)=-\sum_i P(A_i)\log_2 P(A_i)\approx 1.955\ \text{bit},
\]
\[
H(B)=-\sum_j P(B_j)\log_2 P(B_j)\approx 1.560\ \text{bit}.
\]

互信息
\[
I(A;B)=H(A)+H(B)-H(A,B)\approx 1.955+1.560-3.447=0.068\ \text{bit}.
\]

\medskip
\exercise{习题 7.4}
\textbf{题目：}
证明：
\[
I(X;Y)=H(X)-H(X|Y)=H(Y)-H(Y|X)=H(X)+H(Y)-H(X,Y).
\]

\par\textbf{证明：}
由条件熵定义
\[
H(X|Y)=H(X,Y)-H(Y),
\]
因此
\[
H(X)-H(X|Y)=H(X)-[H(X,Y)-H(Y)]=H(X)+H(Y)-H(X,Y).
\]
同理
\[
H(Y)-H(Y|X)=H(Y)-[H(X,Y)-H(X)]=H(X)+H(Y)-H(X,Y).
\]
故三式相等，即
\[
I(X;Y)=H(X)-H(X|Y)=H(Y)-H(Y|X)=H(X)+H(Y)-H(X,Y).
\]

\medskip
\exercise{习题 7.5}
\textbf{题目：}
已知一离散无记忆信源
\[
\{x_1,x_2,x_3\},
\qquad
\{p_1,p_2,p_3\}=\{0.45,0.35,0.20\}.
\]
求：
\begin{enumerate}
\item 信源熵 $H(X)$；
\item 进行哈夫曼编码，并求编码效率。
\end{enumerate}

\par\textbf{解：}
信源熵为
\[
H(X)=-0.45\log_2 0.45-0.35\log_2 0.35-0.20\log_2 0.20\approx 1.51\ \text{bit}.
\]

一种哈夫曼编码可取
\[
x_1\to 1,\qquad x_2\to 01,\qquad x_3\to 00.
\]
平均码长为
\[
\bar L=0.45\times 1+(0.35+0.20)\times 2=1.55\ \text{bit}.
\]
编码效率为
\[
\eta=\frac{H(X)}{\bar L}=\frac{1.51}{1.55}\approx 97.4\%.
\]

\medskip
\exercise{习题 7.6}
\textbf{题目：}
设有一离散无记忆信源
\[
\{x_1,x_2,x_3,x_4,x_5,x_6,x_7\},
\]
概率分别为
\[
0.20,\ 0.19,\ 0.18,\ 0.17,\ 0.15,\ 0.10,\ 0.01.
\]
求：
\begin{enumerate}
\item 信源符号熵 $H(X)$；
\item 哈夫曼编码及编码效率。
\end{enumerate}

\par\textbf{解：}
熵为
\[
\begin{aligned}
H(X)=&-0.20\log_2 0.20-0.19\log_2 0.19-0.18\log_2 0.18 \\
&-0.17\log_2 0.17-0.15\log_2 0.15-0.10\log_2 0.10-0.01\log_2 0.01
\approx 2.61\ \text{bit}.
\end{aligned}
\]

可取一组哈夫曼编码为
\[
x_1\to 11,\quad
x_2\to 10,\quad
x_3\to 011,\quad
x_4\to 010,\quad
x_5\to 001,\quad
x_6\to 0001,\quad
x_7\to 0000.
\]
平均码长
\[
\bar L=(0.20+0.19)\times 2+(0.18+0.17+0.15)\times 3+(0.10+0.01)\times 4=2.72\ \text{bit}.
\]
编码效率
\[
\eta=\frac{2.61}{2.72}\approx 96.0\%.
\]

\medskip
\exercise{习题 7.7}
\textbf{题目：}
试确定能重构信号
\[
x(t)=\sinc(2000t)
\]
所需的最低采样频率。

\par\textbf{解：}
由
\[
\mathcal{F}\{\sinc(2000t)\}=\frac{1}{2000}\rect\!\left(\frac{f}{2000}\right),
\]
可知该信号带宽为
\[
f_m=1000\ \text{Hz}.
\]
按奈奎斯特采样定理，最低采样频率为
\[
f_s=2f_m=2000\ \text{Hz}.
\]

\medskip
\exercise{习题 7.8}
\textbf{题目：}
已知信号
\[
s(t)=10m(t)\cos(2000\pi t),
\]
其中 $m(t)$ 是带宽为 $100$ Hz 的基带信号。对 $s(t)$ 进行带通信号采样，求频谱不交叠时所需的最低采样频率。

\par\textbf{解：}
载波频率为
\[
f_c=1000\ \text{Hz},
\]
因此 $s(t)$ 的频带范围为
\[
900\sim1100\ \text{Hz},
\]
带宽
\[
B=200\ \text{Hz}.
\]
按带通信号采样定理，需找整数 $n$ 使
\[
\frac{2f_H}{n}\le f_s\le \frac{2f_L}{n-1}.
\]
取
\[
f_L=900,\qquad f_H=1100.
\]
当 $n=5$ 时，
\[
\frac{2f_H}{5}=440\ \text{Hz},\qquad \frac{2f_L}{4}=450\ \text{Hz},
\]
存在可行区间，因此最低采样频率为
\[
f_{s,\min}=440\ \text{Hz}.
\]

\medskip
\exercise{习题 7.9}
\textbf{题目：}
已知一个 12 路载波电话占有频率范围
\[
60\sim108\ \text{kHz},
\]
求其最低采样频率。

\par\textbf{解：}
该带通信号频带为
\[
f_L=60\ \text{kHz},\qquad f_H=108\ \text{kHz},
\]
带宽为
\[
B=48\ \text{kHz}.
\]
仍由带通信号采样定理，
\[
\frac{2f_H}{n}\le f_s\le \frac{2f_L}{n-1}.
\]
当 $n=2$ 时，
\[
\frac{2f_H}{2}=108\ \text{kHz},\qquad \frac{2f_L}{1}=120\ \text{kHz},
\]
故可取最小采样频率
\[
f_{s,\min}=108\ \text{kHz}.
\]

\medskip
\exercise{习题 7.10}
\textbf{题目：}
已知正弦信号幅度为 $3.25$ V，将它输入到一个 8 电平均匀量化器。设采样率
\[
f_s=8\ \text{kHz},
\]
正弦信号频率
\[
f=800\ \text{Hz},
\]
试画出量化输出波形。

\par\textbf{解：}
输入信号为
\[
x(t)=3.25\sin(1600\pi t).
\]
采样时刻为
\[
t_i=\frac{i}{f_s}=\frac{i}{8000},
\]
故采样值为
\[
x_i=3.25\sin\!\left(1600\pi\cdot \frac{i}{8000}\right)=3.25\sin\frac{i\pi}{5}.
\]
一个周期内有 10 个采样点，对应采样值依次为
\[
0,\ 1.91,\ 3.09,\ 3.09,\ 1.91,\ 0,\ -1.91,\ -3.09,\ -3.09,\ -1.91.
\]

8 电平均匀量化后，可对应到重构电平
\[
0.5,\ 1.5,\ 3.5,\ 3.5,\ 1.5,\ 0.5,\ -1.5,\ -3.5,\ -3.5,\ -1.5.
\]
因此量化输出波形是按上述 10 个离散电平依次保持的阶梯波。

\medskip
\exercise{习题 7.11}
\textbf{题目：}
已知模拟信号采样值 $X$ 的概率密度为题图 7.11 所示三角分布：
\[
p(x)=1-|x|,\qquad |x|\le 1,
\]
其余为 0。$X$ 通过一个四电平均匀量化器后成为 $Y$。求：
\begin{enumerate}
\item 量化器输入信号功率 $S=E[X^2]$ 与输出信号功率 $S_q=E[Y^2]$；
\item 量化噪声平均功率 $N_q=E[(Y-X)^2]$；
\item 量化信噪比 $S/N_q$（单位：dB）。
\end{enumerate}

\par\textbf{解：}
输入功率
\[
S=E[X^2]=\int_{-1}^{1}x^2p(x)\,dx
=2\int_0^1 x^2(1-x)\,dx
=\frac16.
\]

四电平均匀量化后，输出可表示为
\[
Y=\operatorname{sgn}(X)\cdot \frac14(2+W),
\]
其中
\[
W=\operatorname{sgn}\!\left(|X|-\frac12\right)\in\{\pm1\}.
\]
有
\[
P\!\left(|X|<\frac12\right)=2\int_0^{0.5}(1-x)\,dx=\frac34,
\]
故
\[
E[W]=-\frac34+\frac14=-\frac12,\qquad E[W^2]=1.
\]
于是
\[
S_q=E[Y^2]=\frac1{16}E[(2+W)^2]
=\frac{4+4E[W]+E[W^2]}{16}
=\frac{3}{16}.
\]

量化噪声功率
\[
N_q=E[(Y-X)^2]=E[X^2]-2E[XY]+E[Y^2].
\]
由题解计算可得
\[
E[XY]=\frac16,
\]
故
\[
N_q=\frac16-2\cdot \frac16+\frac{3}{16}=\frac1{48}.
\]

量化信噪比为
\[
\frac{S}{N_q}=\frac{1/6}{1/48}=8,
\]
折合分贝
\[
10\log_{10}8\approx 9.03\ \text{dB}.
\]

\medskip
\exercise{习题 7.12}
\textbf{题目：}
若输入于 A 律 13 折线 PCM 编码器的正弦信号
\[
x(t)=\sin(1600\pi t),
\]
以
\[
f_s=8\ \text{kHz}
\]
的速率采样，得到采样序列
\[
x(n)=\sin(0.2\pi n),\qquad n=0,1,\dots,10.
\]
试求该 PCM 编码器输出的码组序列。

\par\textbf{解：}
对应采样值依次为
\[
0,\ 0.5878,\ 0.9511,\ 0.9511,\ 0.5878,\ 0,\ -0.5878,\ -0.9511,\ -0.9511,\ -0.5878,\ 0.
\]
按 A 律 13 折线 PCM 编码规则查表编码，可得到输出码组序列：
\[
\begin{aligned}
&10000000\ 11110010\ 11111110\ 11111110\ 11110010\ 10000000\\
&01110010\ 01111110\ 01111110\ 01110010\ 00000000
\end{aligned}
\]
共 80 bit。

\medskip
\exercise{习题 7.13}
\textbf{题目：}
某 A 律 13 折线 PCM 编码器的设计输入范围是
\[
[-5,+5]\ \text{V},
\]
若采样值
\[
x=+1.2\ \text{V},
\]
求其编码输出及译码输出值。

\par\textbf{解：}
由于 $x>0$，极性码为
\[
1.
\]
又因为
\[
1.2\in[0.625,1.25],
\]
故段落码为
\[
101.
\]
在该段内再细分 16 小段，$1.2$ 落在倒数第 2 小段，对应段内码
\[
1110.
\]
因此输出码组为
\[
11011110.
\]

译码输出为该量化小段中心值，约为
\[
\hat x\approx 1.191\ \text{V}.
\]

\medskip
\exercise{习题 7.14}
\textbf{题目：}
在 CD 播放机中，假设音乐采用采样率
\[
44.1\ \text{kHz},
\]
每个采样 16 bit 的均匀量化线性编码器进行量化、编码。求 50 分钟音乐所需比特数，并求量化信噪比。

\par\textbf{解：}
50 分钟内所需比特数为
\[
N=44.1\times 10^3\times 16\times 50\times 60
=2.1168\times 10^9\ \text{bit}.
\]

16 bit 均匀量化的量化信噪比近似为
\[
SNR_q\approx 6.02n=6.02\times 16=96.32\ \text{dB},
\]
更精确写为
\[
SNR_q\approx 96.33\ \text{dB}.
\]

\medskip
\exercise{习题 7.15}
\textbf{题目：}
一个离散无记忆信源的字符集为
\[
\{-5,-3,-1,0,1,2,3\},
\]
相应概率为
\[
\{0.08,0.20,0.15,0.03,0.12,0.02,0.40\}.
\]
求：
\begin{enumerate}
\item 信源熵 $H(x)$；
\item 该信源的哈夫曼编码；
\item 平均码长 $R$ 及编码效率 $\eta=H(x)/R$；
\item 若再按量化规则
\[
\hat x=
\begin{cases}
-2,&x=-5,-3,\\
0,&x=-1,0,1,\\
2,&x=2,3,
\end{cases}
\]
求量化后信源熵。
\end{enumerate}

\par\textbf{解：}
信源熵为
\[
\begin{aligned}
H(x)=&-0.08\log_2 0.08-0.20\log_2 0.20-0.15\log_2 0.15-0.03\log_2 0.03\\
&-0.12\log_2 0.12-0.02\log_2 0.02-0.40\log_2 0.40
\approx 2.33\ \text{bit}.
\end{aligned}
\]

一种哈夫曼编码可取
\[
3\to 1,\quad -3\to 011,\quad -1\to 010,\quad 1\to 001,\quad -5\to 0001,\quad 0\to 00001,\quad 2\to 00000.
\]

平均码长为
\[
R=0.40+3(0.20+0.15+0.12)+4\times 0.08+5(0.03+0.02)=2.38\ \text{bit}.
\]
编码效率
\[
\eta=\frac{2.33}{2.38}\approx 97.9\%.
\]

量化后只有 3 个输出值，其概率分别为
\[
P(\hat x=-2)=0.28,\qquad P(\hat x=0)=0.30,\qquad P(\hat x=2)=0.42.
\]
故量化后信源熵为
\[
H(\hat x)=-0.28\log_2 0.28-0.30\log_2 0.30-0.42\log_2 0.42\approx 1.56\ \text{bit}.
\]

\medskip
\exercise{习题 7.16}
\textbf{题目：}
分别对 10 路话音信号（每路最高频率分量为 4 kHz）以奈奎斯特速率采样，采样后均匀量化、线性编码，然后将此 10 路 PCM 信号与两路 180 kbit/s 的数据以时分复用方式复用，时分复用输出送至 BPSK 调制器。若要求 BPSK 信号的功率谱主瓣宽度为
\[
(100\pm 1)\ \text{MHz},
\]
其中 100 MHz 是载频，求每路 PCM 信号允许的最大量化电平数 $M$。

\par\textbf{解：}
主瓣宽度从 $100-1$ 到 $100+1$ MHz，总宽度为
\[
2\ \text{MHz}.
\]
对矩形脉冲 BPSK，
\[
\text{主瓣宽度}=2R_b,
\]
故时分复用后的总码率为
\[
R_b=1\ \text{Mbit/s}=1000\ \text{kbit/s}.
\]

每路语音奈奎斯特采样率为
\[
8\ \text{kHz}.
\]
若每个样值编码为 $\log_2 M$ bit，则每路 PCM 速率为
\[
8\log_2 M\ \text{kbit/s}.
\]
10 路 PCM 总速率再加上两路 180 kbit/s 数据，可得
\[
10\times 8\log_2 M+2\times 180=1000.
\]
即
\[
80\log_2 M=640,
\qquad
\log_2 M=8.
\]
因此
\[
M=256.
\]

\medskip
\exercise{习题 7.17}
\textbf{题目：}
对 10 路模拟信号分别进行 A 律 13 折线 PCM，然后进行时分复用，再以二进制方式经滚降因子为 0.5 的升余弦频谱滚降系统进行无码间干扰传输。该升余弦基带系统的截止频率为
\[
480\ \text{kHz}.
\]
求：
\begin{enumerate}
\item 该系统的最大信息传输速率；
\item 允许每路模拟信号的最高频率分量 $f_H$。
\end{enumerate}

\par\textbf{解：}
升余弦系统截止频率满足
\[
\frac{R_b}{2}(1+\alpha)=480\ \text{kHz}.
\]
代入
\[
\alpha=0.5
\]
得
\[
R_b=\frac{480\times 2}{1+0.5}=640\ \text{kbit/s}.
\]

又由于 10 路 A 律 13 折线 PCM，每个采样值编码为 8 bit，因此总码率为
\[
10\times 2f_H\times 8.
\]
令其等于 640 kbit/s：
\[
10\times 2f_H\times 8=640,
\]
解得
\[
f_H=\frac{640}{160}=4\ \text{kHz}.
\]

\medskip
